\section{Character minors of imaginary quadratic class groups}
\label{sec:class-group-minors}
This section separates an algebraic obstruction to nonvanishing from
the analytic question of finding prime ideals in prescribed classes.
The distinction is already decisive for the ordinary ideal class group.

\begin{proposition}\label{prop:abelian-all-minors}
Let $A$ be a finite abelian group. Every square minor of its complex
character table is nonzero if and only if $|A|$ is one or a prime.
For a subset of distinct columns in a prime-order character table,
all square minors are nonzero without an analytic hypothesis.
\end{proposition}
\begin{proof}
If $|A|$ is composite, choose a proper nontrivial subgroup $H$, an
element $h\ne1$ in $H$, and a nontrivial character of $A/H$.
Pulling that character back gives $\chi\ne1$ with $\chi(h)=1$.
The rows $1,\chi$ and columns $1,h$ give the singular matrix
$\left(\begin{smallmatrix}1&1\\1&1\end{smallmatrix}\right)$.
If $|A|$ is prime, the table is the prime Fourier matrix, up to row
and column order. Chebotarev's theorem gives the conclusion; see
\citet[Lemma~1.3, author preprint pp.~3--4]{Tao2005}.
The order-one case is immediate.
\end{proof}

Here is an explicit counterexample in which both columns come from
distinct classes of unramified degree-one prime ideals.
Put $K=\mathbb Q(\sqrt{-14})$, whose integer ring is
$\mathbb Z[\sqrt{-14}]$, and write $s=\sqrt{-14}$.
The prime ideals $\mathfrak p_2=(2,s)$ and
$\mathfrak p_3=(3,s-1)$ have norms two and three.
Unique ideal factorization gives
\[
 \mathfrak p_2^2=(2),\qquad
 (2+s)=\mathfrak p_2\mathfrak p_3^2.
\]
For the second equality, the norm is $18$, the indicated primes
divide $2+s$, and its conjugate prime above three does not divide it.
The class of $\mathfrak p_2$ is nontrivial, since
$a^2+14b^2=2$ has no integer solution. It follows that
$[\mathfrak p_3]$ has order four, and that
$H=\langle[\mathfrak p_2]\rangle$ is a proper subgroup of the class group.
Choose any nontrivial character $\chi$ of the quotient by $H$.

The ideal $(3+s)$ has norm $23$ and is a principal prime.
The element $4+3s$ has norm $142=2\cdot71$, so
\[
 (4+3s)=\mathfrak p_2\mathfrak p_{71}
\]
for a prime ideal $\mathfrak p_{71}$ of norm $71$.
Thus $[\mathfrak p_{71}]=[\mathfrak p_2]\ne1$.
The two primes of norms $23$ and $71$ are unramified, have residue
degree one, and give distinct columns on which $1$ and $\chi$ agree.
Their $2\times2$ minor vanishes. It remains a minor of every table
formed using a larger cutoff $x\ge71$.
Consequently no eventual cutoff, polynomial in $\log|\operatorname{disc}K|$
or otherwise, can imply the proposed all-minors nonvanishing in general.
This counterexample also treats the interpretation in which repeated
prime ideals from the same class are first removed.

\subsection{What the analytic hypothesis would concern}
Let $\mathfrak m$ be a ray modulus and $E/K$ its ray class field.
The relevant generalized Riemann hypothesis (GRH), denoted
$\mathrm{GRH}(K,\mathfrak m)$, is that every primitive Hecke
$L_K(s,\chi)$ of finite order and
 conductor dividing $\mathfrak m$ has all nontrivial zeros on
 $\operatorname{Re}s=\tfrac12$.
The trivial character includes $\zeta_K$. Artin reciprocity identifies
these characters with those of $\Gamma=\operatorname{Gal}(E/K)$, and
\[
 \zeta_E(s)=\prod_{\chi\in\widehat\Gamma}L_K(s,\chi),
\]
with the corresponding primitive Hecke factors understood.
The Euler-factor identity can be checked directly. At a prime of $K$
with inertia subgroup $I\subseteq\Gamma$ and residue degree $f$,
only characters trivial on $I$ contribute. Their values on Frobenius
are the $f$th roots of unity, each repeated $|\Gamma|/(|I|f)$ times.
Their product is therefore
$(1-(N\mathfrak p)^{-fs})^{-|\Gamma|/(|I|f)}$, precisely the local
factor of $\zeta_E$. The nontrivial primitive Hecke factors are
holomorphic in the open critical strip, and the trivial factor's pole
is at $s=1$, so zeros in that strip cannot cancel in this product.
Thus the named hypothesis may equivalently be imposed on $\zeta_E$.
An effective Chebotarev theorem under this Dedekind-GRH hypothesis is
given by \citet[Theorem~1.1 and Corollary~1.2, corrected author version~3,
pp.~2--3]{GrenieMolteni2019}.

The scope of a Dirichlet-GRH input is explicit in the norm-induced
case. Let $\eta_K$ be the quadratic Dirichlet character of $K$, let
$\psi$ be a Dirichlet character, and let $\Sigma$ contain the rational
primes dividing their conductors or the discriminant of $K$.
Deleting the Euler factors above $\Sigma$ gives
\[
 L_K^{(\Sigma)}(s,\psi\circ N_{K/\mathbb Q})
 =L_{\mathbb Q}^{(\Sigma)}(s,\psi)
  L_{\mathbb Q}^{(\Sigma)}(s,\psi\eta_K).
\]
Indeed, a split prime $l\notin\Sigma$ gives two factors
$(1-\psi(l)l^{-s})^{-1}$ on either side. An inert prime gives
$(1-\psi(l)^2l^{-2s})^{-1}$ on the left; on the right the two
factors have opposite signs because $\eta_K(l)=-1$, giving the
same expression. The omitted finite Euler factors have zeros and
poles only on $\operatorname{Re}s=0$ and hence do not affect GRH in
the open critical strip. Dirichlet GRH for $\psi$ and $\psi\eta_K$
therefore covers this norm-induced character. No such reduction is
asserted for an arbitrary ray-class character of $K$; those characters
remain within the explicitly named additional hypothesis above.

Prime-class coverage, small generators, and nonvanishing of all
character minors are three different statements. Effective
Chebotarev concerns the first; its natural field invariants include
the ray conductor and the discriminant and degree of $E$.
Omitting an unrestricted ray conductor from a uniform bound is an
additional issue. None of these analytic estimates removes the
singular minor proved above. In prime class-group order, on the other
hand, distinct-column nonvanishing holds before any prime bound is used.
